\documentclass[11pt]{article}

%%%%%%%%%%%%%%Include Packages%%%%%%%%%%%%%%%%%%%%%%%%%%
\usepackage{xcolor}
\usepackage{mathtools}
\usepackage[a4paper, total={6in, 8in}, margin=1.25in]{geometry}
\usepackage{amsmath}
\usepackage{amssymb}
\usepackage{paralist}
\usepackage{rsfso}
\usepackage{amsthm}
\usepackage{wasysym}
\usepackage[inline]{enumitem}   
\usepackage{hyperref}
\usepackage{tocloft}
\usepackage{titlesec}
\usepackage{colortbl}
%%%%%%%%%%%%%%%%%%%%%%%%%%%%%%%%%%%%%%%%%%%%%%%%%%%%%%%%



\begin{document}
\Large \noindent \textbf{Physics 5C Midterm 1 Review Session}\\
\normalsize
Department of Physics and Astronomy, UCLA\\
%\noindent TA: Jinyan Miao (jnmiu@ucla.edu)\\
%Office hours: Thursday 10-11 AM $\&$ 1-2 PM at PAB 2-429\\
%Tutoring center hours: Friday 2-3 PM at PAB 2-429 (Zoom room also available)\\
\noindent\rule{8cm}{0.4pt}\\

\noindent Consider the following configuration:\\
A charge $q_1=-1\,$C is fixed at $\vec{x}_1 =(0,0)\,$m.\\
A charge $q_2=1\,$C is fixed at $\vec{x}_2 =(1,-1)\,$m.\\
Jinyan is trying to bring in a test charge $q_3=1\,$C from infinitely far away at $\vec{x}_i = (\infty,0)$\,m to the destination point $\vec{x}_f = (1,0)\,$m. We will discuss this process from two perspectives, first using electric potential and energy, then using electric field and forces.\\

\noindent First we consider the picture of electric potential.\\
\noindent\textbf{Part (a)} Compute the electric potential at $\vec{x}_i$ and $\vec{x}_f$ to the source charges $q_1$ and $q_2$.\\
\noindent\textbf{Part (b)} Compute the electric potential energy stored in $q_3$ at $\vec{x}_i$ and $\vec{x}_f$.\\
\noindent\textbf{Part (c)} How much energy is needed to bring $q_3$ from $\vec{x}_i$ to $\vec{x}_f$? \\
\noindent\textbf{Part (d)} If $q_3$ is fixed at $\vec{x}_f = (1,0)\,$m, compute the potential energy stored in $q_2$. \\
\noindent\textbf{Part (e)} What is the total energy needed to assemble the three-charge system? \footnote{That is, this should be a three-step calculation: Step 1 - Starting with a vacuum, nothing in space, find the energy to put $q_1$ at the origin $\vec{x}_1$; Step 2 - With only $q_1$ fixed at the origin, find the energy needed to bring in $q_2$ from infinity to $\vec{x}_2 = (1,-1)\,$m; Step 3 - With $q_1$ and $q_2$ fixed at their locations, find the energy needed to bring in $q_3$ from infinity to $\vec{x}_f$. Sum these energy to find the total energy stored in such a system.}\\

\noindent Now let us consider the picture of electric field.\\
\noindent\textbf{Part (f)} Draw the electric field lines due to the source charges $q_1$ and $q_2$ combined.\\
\noindent\textbf{Part (g)} Compute the net electric field due to $q_1$ and $q_2$ combined, at $\vec{x}_i$ and at $\vec{x}_f$.\\
\noindent\textbf{Part (h)} Compute the net force that $q_3$ experiences when it is at $\vec{x}_i$ and at $\vec{x}_f$. \\\
\noindent\textbf{Part (i)} Jinyan flips the sign of $q_3$ when it is at $\vec{x}_f$. What gets changed in this picture?\\

\noindent Now we explore the relation between electric field and electric potential. \\
Consider only the source charges $q_1$ and $q_2$, that is, neglect the test charge $q_3$.\\
\noindent\textbf{Part (j)} Find the location (other than very far away) where we have the ``net" electric potential being zero. At that location, find the magnitude of the electric field. Qualitatively describe the relation between electric field and electric potential. 

\newpage
\section*{Solution}

\end{document}


